\uuid{lgCD}
\titre{Primitive}
\auteur{bodin ; croizat ; sacre}
\organisation{exo7}
\licence{CC BY-NC-SA 4.0}
\datecreate{2022-03-08}
\niveau{L1}
\module{Analyse}
\chapitre{Calcul d'intégrales}
\sousChapitre{Primitives diverses}
\difficulte{2}
\theme{Primitive}
\exo7qcmid{hilisit.equations-differentielles.007}

\contenu{
\begin{question}
\qclassification{127.11}
Soit $F$ une primitive d'une fonction $f$ et $G$ une primitive d'une fonction $g$ sur un intervalle $I$.
Quelles sont les affirmations vraies ?
\begin{answers}
  \bad{Si $f=g$ alors $F=G$.}
  \good{Si $F=G$ alors $f=g$.}
  \bad{Si $f=g^2$ alors $F=G^2$.}
  \good{Si $F=G+C$ (où $C$ est une constante) alors $f=g$.}
\end{answers}
\begin{explanations}
Si $f=g$ alors $F=G+C$ (où $C$ est une constante).
On rappelle que $F'=f$ et $G'=g$, donc si $F=G+C$ alors en dérivant l'égalité on obtient $F'=f = (G+C)'=G'+0=g$. Remarquez par ailleurs que les primitives de $x^2$ sont $\frac{x^3}{3} + C$ (où $C$ est une constante) : ce ne sont les carrés des primitives de $x$ (qui sont $\frac{x^2}{2} + \widetilde{C}$, où $\widetilde{C}$ est une constante).
\end{explanations}
\end{question}
}
